% ═══════════════════════════════════════════════════════════════════════════
% SEQUENZÜBERSICHT - 10 Doppelstunden zur Beobachtung Feb 2026
% Spanisch Spätbeginner (Kl. 10, 11, 12) - Jahrgangsübergreifend
% ERWEITERTE VERSION mit detaillierten Stundenbeschreibungen
% ═══════════════════════════════════════════════════════════════════════════

\documentclass[11pt, a4paper]{article}

\usepackage[utf8]{inputenc}
\usepackage[T1]{fontenc}
\usepackage[ngerman]{babel}
\usepackage{helvet}
\renewcommand{\familydefault}{\sfdefault}

\usepackage[margin=1.5cm, top=1.5cm, bottom=1.8cm]{geometry}
\usepackage[table]{xcolor}
\usepackage{array}
\usepackage{tabularx}
\usepackage{longtable}
\usepackage{booktabs}
\usepackage{tikz}
\usepackage{amssymb}
\usepackage{enumitem}
\usepackage{fancyhdr}
\usepackage{lastpage}
\usepackage{tcolorbox}
\usepackage{multirow}
\usepackage{colortbl}

\tcbuselibrary{skins,breakable}

% Farben
\definecolor{spanisch}{HTML}{c41e3a}
\definecolor{grau}{HTML}{666666}
\definecolor{hellgrau}{HTML}{f5f5f5}
\definecolor{gruppeA}{HTML}{2a7a4b}
\definecolor{gruppeB}{HTML}{1565c0}
\definecolor{puffer}{HTML}{6a0dad}
\definecolor{assessment}{HTML}{b35c00}
\definecolor{beobachtung}{HTML}{c41e3a}
\definecolor{success}{HTML}{2a7a4b}

% Spanisch-Befehl
\newcommand{\es}[1]{\textcolor{spanisch}{\textit{#1}}}

% Info-Boxen
\newtcolorbox{dsbox}[1][]{
    colback=white,
    colframe=spanisch,
    boxrule=0.8pt,
    arc=0pt,
    left=6pt, right=6pt, top=4pt, bottom=4pt,
    fonttitle=\bfseries,
    #1
}

\newtcolorbox{assessmentbox}[1][]{
    colback=yellow!5,
    colframe=assessment,
    boxrule=1.2pt,
    arc=0pt,
    left=6pt, right=6pt, top=4pt, bottom=4pt,
    fonttitle=\bfseries\color{assessment},
    #1
}

\newtcolorbox{pufferbox}[1][]{
    colback=purple!5,
    colframe=puffer,
    boxrule=1.2pt,
    arc=0pt,
    left=6pt, right=6pt, top=4pt, bottom=4pt,
    fonttitle=\bfseries\color{puffer},
    #1
}

\newtcolorbox{beobachtungsbox}[1][]{
    colback=red!5,
    colframe=beobachtung,
    boxrule=1.5pt,
    arc=0pt,
    left=6pt, right=6pt, top=4pt, bottom=4pt,
    fonttitle=\bfseries\color{beobachtung},
    #1
}

% Kopf- und Fußzeile
\pagestyle{fancy}
\fancyhf{}
\renewcommand{\headrulewidth}{0.4pt}
\lhead{\small Sequenzübersicht Spanisch Spätbeginner}
\rhead{\small Beobachtung Februar 2026}
\cfoot{\small Seite \thepage\ von \pageref{LastPage}}

\setlength{\parindent}{0pt}
\setlength{\parskip}{0.3em}

\begin{document}

% ═══════════════════════════════════════════════════════════════════════════
% TITELSEITE
% ═══════════════════════════════════════════════════════════════════════════

\begin{center}
{\color{spanisch}\rule{\textwidth}{2pt}}\\[0.6cm]
{\Huge\bfseries SEQUENZÜBERSICHT}\\[0.2cm]
{\Large Spanisch Spätbeginner -- 10 Doppelstunden}\\[0.4cm]
{\large Vorbereitung Unterrichtsbeobachtung}\\[0.15cm]
{\large\bfseries 3. Februar 2026}\\[0.4cm]
{\color{spanisch}\rule{\textwidth}{2pt}}
\end{center}

\vspace{0.4cm}

% ═══════════════════════════════════════════════════════════════════════════
% KURSINFO
% ═══════════════════════════════════════════════════════════════════════════

\begin{center}
\begin{tabular}{|p{3.5cm}|p{11.5cm}|}
\hline
\rowcolor{hellgrau}
\textbf{Kurs} & Spanisch Spätbeginner (jahrgangsübergreifend) \\
\hline
\textbf{Klassenstufen} & 10, 11, 12 \\
\hline
\textbf{Schülerzahl} & 10 (1m, 9w) \\
\hline
\textbf{Lehrwerk} & A\_tope.com (Cornelsen) \\
\hline
\textbf{Wochenstunden} & 4 (Mo + Di je Doppelstunde) \\
\hline
\rowcolor{hellgrau}
\textbf{Differenzierung} & \textcolor{gruppeA}{\textbf{Gruppe A}} (7 SuS): Lernjahr 1, A2-Niveau \quad \textcolor{gruppeB}{\textbf{Gruppe B}} (3 SuS): Lernjahr 2--3, B1-Niveau \\
\hline
\end{tabular}
\end{center}

\vspace{0.4cm}

% ═══════════════════════════════════════════════════════════════════════════
% TIMELINE-VISUALISIERUNG
% ═══════════════════════════════════════════════════════════════════════════

{\large\bfseries\color{spanisch} Timeline: 5 Wochen bis zur Beobachtung}

\vspace{0.3cm}

\begin{center}
\begin{tikzpicture}[scale=0.88, transform shape]
% Hintergrund-Rechtecke für Wochen
\fill[hellgrau!50] (-0.5,-0.8) rectangle (3.5,1.2);
\fill[hellgrau!50] (3.7,-0.8) rectangle (7.7,1.2);
\fill[hellgrau!50] (7.9,-0.8) rectangle (11.9,1.2);
\fill[hellgrau!50] (12.1,-0.8) rectangle (16.1,1.2);
\fill[red!10] (16.3,-0.8) rectangle (20.3,1.2);

% Wochen-Header
\node[font=\small\bfseries] at (1.5,0.9) {Woche 1};
\node[font=\small\bfseries] at (5.7,0.9) {Woche 2};
\node[font=\small\bfseries] at (9.9,0.9) {Woche 3};
\node[font=\small\bfseries] at (14.1,0.9) {Woche 4};
\node[font=\small\bfseries, color=beobachtung] at (18.3,0.9) {Woche 5};

% Datums-Zeile
\node[font=\tiny, color=grau] at (1.5,0.5) {05.--06.01.};
\node[font=\tiny, color=grau] at (5.7,0.5) {12.--13.01.};
\node[font=\tiny, color=grau] at (9.9,0.5) {19.--20.01.};
\node[font=\tiny, color=grau] at (14.1,0.5) {26.--27.01.};
\node[font=\tiny, color=grau] at (18.3,0.5) {02.--03.02.};

% Haupt-Timeline
\draw[thick, spanisch] (0,0) -- (20,0);

% Doppelstunden-Marker
% W1: DS 1+2
\fill[grau] (0.5,0) circle (5pt); \node[font=\scriptsize, above] at (0.5,0.15) {DS1};
\node[font=\tiny, below] at (0.5,-0.15) {Mo};
\fill[grau] (2.5,0) circle (5pt); \node[font=\scriptsize, above] at (2.5,0.15) {DS2};
\node[font=\tiny, below] at (2.5,-0.15) {Di};

% W2: DS 3+4
\fill[grau] (4.7,0) circle (5pt); \node[font=\scriptsize, above] at (4.7,0.15) {DS3};
\node[font=\tiny, below] at (4.7,-0.15) {Mo};
\fill[grau] (6.7,0) circle (5pt); \node[font=\scriptsize, above] at (6.7,0.15) {DS4};
\node[font=\tiny, below] at (6.7,-0.15) {Di};

% W3: DS 5 (Assessment) + DS 6
\fill[assessment] (8.9,0) circle (6pt); \node[font=\scriptsize\bfseries, above, color=assessment] at (8.9,0.15) {DS5};
\node[font=\tiny, below] at (8.9,-0.15) {Mo};
\fill[grau] (10.9,0) circle (5pt); \node[font=\scriptsize, above] at (10.9,0.15) {DS6};
\node[font=\tiny, below] at (10.9,-0.15) {Di};

% W4: DS 7+8 (Puffer)
\fill[puffer] (13.1,0) circle (5pt); \node[font=\scriptsize, above, color=puffer] at (13.1,0.15) {DS7};
\node[font=\tiny, below] at (13.1,-0.15) {Mo};
\fill[puffer] (15.1,0) circle (5pt); \node[font=\scriptsize, above, color=puffer] at (15.1,0.15) {DS8};
\node[font=\tiny, below] at (15.1,-0.15) {Di};

% W5: DS 9 + DS 10 (Beobachtung)
\fill[grau] (17.3,0) circle (5pt); \node[font=\scriptsize, above] at (17.3,0.15) {DS9};
\node[font=\tiny, below] at (17.3,-0.15) {Mo};
\fill[beobachtung] (19.3,0) circle (7pt); \node[font=\scriptsize\bfseries, above, color=beobachtung] at (19.3,0.15) {DS10};
\node[font=\tiny, below] at (19.3,-0.15) {Di};

% Legende
\node[font=\tiny, right] at (0,-0.6) {
    \textcolor{grau}{$\bullet$} Bucharbeit/Brücke \quad
    \textcolor{assessment}{$\bullet$} \textbf{Assessment} (45', 25 BE) \quad
    \textcolor{puffer}{$\bullet$} Puffer \quad
    \textcolor{beobachtung}{$\bullet$} \textbf{Beobachtung}
};
\end{tikzpicture}
\end{center}

\vspace{0.4cm}

% ═══════════════════════════════════════════════════════════════════════════
% KOMPAKTÜBERSICHT (TABELLE)
% ═══════════════════════════════════════════════════════════════════════════

{\large\bfseries\color{spanisch} Übersicht aller 10 Doppelstunden}

\vspace{0.2cm}

\renewcommand{\arraystretch}{1.3}
\small
\begin{longtable}{|c|p{2cm}|p{3.2cm}|p{4cm}|p{4.2cm}|}
\hline
\rowcolor{spanisch}
\textcolor{white}{\textbf{DS}} &
\textcolor{white}{\textbf{Datum}} &
\textcolor{white}{\textbf{Thema}} &
\textcolor{white}{\textbf{Lernziele (Kurzform)}} &
\textcolor{white}{\textbf{Besonderheiten}} \\
\hline
\endfirsthead
\hline
\rowcolor{spanisch}
\textcolor{white}{\textbf{DS}} &
\textcolor{white}{\textbf{Datum}} &
\textcolor{white}{\textbf{Thema}} &
\textcolor{white}{\textbf{Lernziele (Kurzform)}} &
\textcolor{white}{\textbf{Besonderheiten}} \\
\hline
\endhead

\textbf{1} & 05.01. (Mo) & Bucharbeit parallel &
A: \es{Mi barrio}, \es{hay}+Ort; B: Berufe, \es{ser/estar} &
Erste Stunde nach Ferien \\
\hline

\textbf{2} & 06.01. (Di) & Fortführung + Übergang &
Lektionen festigen, erstes gemeinsames Element &
B als Experten einsetzen \\
\hline

\textbf{3} & 12.01. (Mo) & Bucharbeit + Ankündigung &
Lektionen abschließen, Assessment-Kriterien kennen &
\textcolor{assessment}{Assessment-Info!} \\
\hline

\textbf{4} & 13.01. (Di) & Letzte Bucharbeit &
A: \es{hay/ser/estar} sichern; B: Vortrag strukturieren &
\textcolor{assessment}{Letzte Übung!} \\
\hline

\rowcolor{yellow!15}
\textbf{5} & 19.01. (Mo) & \textbf{ASSESSMENT} &
A: Dialog; B: Kurzvortrag &
\textcolor{assessment}{\textbf{45', 25 BE, 15 NP}} \\
\hline

\textbf{6} & 20.01. (Di) & Brückenthema &
\es{ser} vs. \es{estar} + Adj., Personen beschreiben &
\textcolor{gruppeB}{B-SuS als Tutoren!} \\
\hline

\rowcolor{purple!8}
\textbf{7} & 26.01. (Mo) & \textcolor{puffer}{Cultura y Diversión} &
Kulturelles Wissen, Wortschatz anwenden &
\textcolor{puffer}{\textbf{PUFFER!}} \\
\hline

\rowcolor{purple!8}
\textbf{8} & 27.01. (Di) & \textcolor{puffer}{Creatividad y Juego} &
\es{Mi persona favorita}, Schreibkompetenz &
\textcolor{puffer}{\textbf{PUFFER!}} \\
\hline

\textbf{9} & 02.02. (Mo) & Vocabulario y Gramática &
Vokabel-Bingo, Grammatik-Stationen, Schreibübung &
Spiele + Kurze Ankündigung \\
\hline

\rowcolor{red!12}
\textbf{10} & 03.02. (Di) & \textbf{BEOBACHTUNG} &
Kommunikative Kompetenz, 3. Person Singular &
\textcolor{beobachtung}{\textbf{EXTERNE BEOBACHTUNG!}} \\
\hline

\end{longtable}

\normalsize
\renewcommand{\arraystretch}{1.0}

\newpage

% ═══════════════════════════════════════════════════════════════════════════
% DETAILLIERTE STUNDENBESCHREIBUNGEN
% ═══════════════════════════════════════════════════════════════════════════

{\Large\bfseries\color{spanisch} Detaillierte Stundenbeschreibungen}

\vspace{0.4cm}

% ═══════════════════════════════════════════════════════════════════════════
% PHASE 1: BUCHARBEIT & ASSESSMENT
% ═══════════════════════════════════════════════════════════════════════════

\textbf{\Large\color{spanisch} PHASE 1: Bucharbeit \& Assessment (DS 1--5)}

\vspace{0.3cm}

% ---------- DS 1 ----------
\begin{dsbox}[title={DS 1 -- Bucharbeit parallel (Mo, 05.01.2026)}]

\textbf{Thema:} Parallele Bucharbeit -- A: \es{Mi barrio} (S. 30--31) | B: Berufe (S. 102--103)

\vspace{0.3cm}

\textbf{LERNZIELE:}
\begin{itemize}[leftmargin=1.5em, nosep]
    \item[\textcolor{gruppeA}{A:}] SuS können Orte im Stadtviertel benennen und \es{hay} + Ortsangabe verwenden
    \item[\textcolor{gruppeA}{A:}] SuS können \es{vuestro/a/os/as} korrekt mit Bezugswort kongruent bilden
    \item[\textcolor{gruppeB}{B:}] SuS können Berufe und notwendige Eigenschaften auf Spanisch beschreiben
    \item[\textcolor{gruppeB}{B:}] SuS können \es{ser/estar} im Kontext Berufsbeschreibung korrekt anwenden
\end{itemize}

\vspace{0.2cm}

\textbf{VERLAUF (90 Min):}

\begin{tabular}{|p{1.2cm}|p{5.5cm}|p{5.5cm}|p{1.5cm}|}
\hline
\rowcolor{hellgrau}
\textbf{Zeit} & \textbf{Gruppe A} & \textbf{Gruppe B} & \textbf{L bei} \\
\hline
0--5 & \multicolumn{2}{l|}{Gemeinsame Begrüßung, Tagesübersicht, Gruppenaufteilung} & Alle \\
\hline
5--20 & Aufg. 9: \es{vuestro}-Fragen formulieren & Aufg. 6a: Schulsystem erklären & A $\rightarrow$ B \\
\hline
20--35 & Aufg. 10: Partnerdialog Madrid (S.136) & Aufg. 6b: Vergleich DE/ES & B $\rightarrow$ A \\
\hline
35--45 & \multicolumn{2}{l|}{\textit{Pause (5 Min)}} & -- \\
\hline
45--60 & Dialoge präsentieren, ser/estar-Korrektur & Mediación (Aufg. 7) einführen & A $\rightarrow$ B \\
\hline
60--80 & Aufg. 11: Blogeintrag \es{Mi barrio} & Mediación abschließen & B $\rightarrow$ A \\
\hline
80--90 & \multicolumn{2}{l|}{Gemeinsamer Abschluss, HA-Ankündigung} & Alle \\
\hline
\end{tabular}

\vspace{0.2cm}

\textbf{AKTIVITÄTEN:}
\begin{itemize}[leftmargin=1.5em, nosep]
    \item Lehrwerksarbeit mit Selbstkontrolle (Lösungen im Anhang)
    \item Partnerdialoge mit Rollenkarten (A: S.136)
    \item Schriftliche Produktion (Blogeintrag/Mediación)
\end{itemize}

\vspace{0.2cm}

\textbf{INTENTIONEN:}
\begin{itemize}[leftmargin=1.5em, nosep]
    \item Reaktivierung nach den Ferien durch vertraute Bucharbeit
    \item Parallele Differenzierung ermöglicht niveaugerechtes Arbeiten
    \item L rotiert zwischen Gruppen (Timer: alle 10--15 Min wechseln)
\end{itemize}

\end{dsbox}

\vspace{0.4cm}

% ---------- DS 2 ----------
\begin{dsbox}[title={DS 2 -- Fortführung + Übergang (Di, 06.01.2026)}]

\textbf{Thema:} Weiterführung Bucharbeit + Erstes gemeinsames Element

\vspace{0.3cm}

\textbf{LERNZIELE:}
\begin{itemize}[leftmargin=1.5em, nosep]
    \item[\textcolor{gruppeA}{A:}] SuS können Hörtexte zu \es{Mi barrio} verstehen und Informationen entnehmen
    \item[\textcolor{gruppeB}{B:}] SuS können ihre Mediación-Ergebnisse präsentieren
    \item[Alle:] SuS können grundlegende Adjektive mit \es{ser} verwenden
\end{itemize}

\vspace{0.2cm}

\textbf{VERLAUF (90 Min):}

\begin{tabular}{|p{1.2cm}|p{5.5cm}|p{5.5cm}|p{1.5cm}|}
\hline
\rowcolor{hellgrau}
\textbf{Zeit} & \textbf{Gruppe A} & \textbf{Gruppe B} & \textbf{L bei} \\
\hline
0--5 & \multicolumn{2}{l|}{Begrüßung, HA-Kontrolle (Stichprobe)} & Alle \\
\hline
5--20 & Hörübung S. 31, Aufg. 12--13 & Mediación-Präsentation vorbereiten & A $\rightarrow$ B \\
\hline
20--35 & Hörübung auswerten & Mediación vorstellen (B an A) & Alle \\
\hline
35--45 & \multicolumn{2}{l|}{\textit{Pause (5 Min) + Umgruppierung}} & -- \\
\hline
45--65 & \multicolumn{2}{l|}{\textbf{GEMEINSAM:} Vokabelspiel (Kahoot/Quizlet) -- Orte + Eigenschaften} & Alle \\
\hline
65--80 & \multicolumn{2}{l|}{\textbf{GEMEINSAM:} \es{ser} + Adjektiv einführen: \es{¿Cómo eres?}} & Alle \\
\hline
80--90 & \multicolumn{2}{l|}{Abschluss, HA (Vokabeln lernen)} & Alle \\
\hline
\end{tabular}

\vspace{0.2cm}

\textbf{AKTIVITÄTEN:}
\begin{itemize}[leftmargin=1.5em, nosep]
    \item Hörverstehen mit Arbeitsblatt
    \item Gemeinsames Vokabelspiel (digital oder analog)
    \item B-SuS als Experten: Erklären Mediación-Konzept
\end{itemize}

\vspace{0.2cm}

\textbf{INTENTIONEN:}
\begin{itemize}[leftmargin=1.5em, nosep]
    \item Erstes gemeinsames Element schafft Klassengemeinschaft
    \item B-SuS als Experten stärkt Selbstbewusstsein und gibt A Sprachmodelle
    \item Vorbereitung auf gemeinsames Brückenthema (DS 6)
\end{itemize}

\end{dsbox}

\vspace{0.4cm}

% ---------- DS 3 ----------
\begin{dsbox}[title={DS 3 -- Bucharbeit + Assessment-Ankündigung (Mo, 12.01.2026)}]

\textbf{Thema:} Lektionen abschließen + Assessment-Vorbereitung

\vspace{0.3cm}

\textbf{LERNZIELE:}
\begin{itemize}[leftmargin=1.5em, nosep]
    \item[\textcolor{gruppeA}{A:}] SuS können \es{ser/estar/hay} situationsgerecht unterscheiden
    \item[\textcolor{gruppeB}{B:}] SuS können Berufswunsch mit Begründung präsentieren
    \item[Alle:] SuS kennen die Bewertungskriterien für das Assessment
\end{itemize}

\vspace{0.2cm}

\textbf{VERLAUF (90 Min):}

\begin{tabular}{|p{1.2cm}|p{5.5cm}|p{5.5cm}|p{1.5cm}|}
\hline
\rowcolor{hellgrau}
\textbf{Zeit} & \textbf{Gruppe A} & \textbf{Gruppe B} & \textbf{L bei} \\
\hline
0--5 & \multicolumn{2}{l|}{Begrüßung, Überblick} & Alle \\
\hline
5--25 & Übungen ser/estar/hay (Differenzierung) & Vortragsgliederung erstellen & A $\rightarrow$ B \\
\hline
25--40 & Partnerdialog-Übung (Assessment-Format) & Vortrag üben (Peer-Feedback) & B $\rightarrow$ A \\
\hline
40--50 & \multicolumn{2}{l|}{\textit{Pause (5 Min) + Umgruppierung}} & -- \\
\hline
50--70 & \multicolumn{2}{l|}{\textbf{ASSESSMENT-ANKÜNDIGUNG:} Formate, Kriterien, Notenskala erklären} & Alle \\
\hline
70--85 & \multicolumn{2}{l|}{Fragen klären, Bewertungsbögen ausgeben, Übungshinweise} & Alle \\
\hline
85--90 & \multicolumn{2}{l|}{Abschluss, HA: Vorbereitung üben} & Alle \\
\hline
\end{tabular}

\vspace{0.2cm}

\textbf{AKTIVITÄTEN:}
\begin{itemize}[leftmargin=1.5em, nosep]
    \item Grammatik-Übungen (differenziert nach Niveau)
    \item Assessment-Formate vorstellen: A = Partnerprüfung, B = Kurzvortrag
    \item Bewertungskriterien transparent machen (Bögen ausgeben)
\end{itemize}

\vspace{0.2cm}

\textbf{INTENTIONEN:}
\begin{itemize}[leftmargin=1.5em, nosep]
    \item Transparenz schafft Sicherheit und reduziert Prüfungsangst
    \item Eine Woche Vorbereitungszeit ist angemessen
    \item SuS können gezielt üben, wenn Kriterien bekannt sind
\end{itemize}

\end{dsbox}

\vspace{0.4cm}

% ---------- DS 4 ----------
\begin{dsbox}[title={DS 4 -- Letzte Bucharbeit vor Assessment (Di, 13.01.2026)}]

\textbf{Thema:} Intensive Übungsstunde -- Generalprobe

\vspace{0.3cm}

\textbf{LERNZIELE:}
\begin{itemize}[leftmargin=1.5em, nosep]
    \item[\textcolor{gruppeA}{A:}] SuS können einen 3-minütigen Dialog über \es{Mi barrio} führen
    \item[\textcolor{gruppeB}{B:}] SuS können einen 3-4-minütigen Kurzvortrag über \es{Mi profesión ideal} halten
    \item[Alle:] SuS können ihre eigene Leistung anhand der Kriterien einschätzen
\end{itemize}

\vspace{0.2cm}

\textbf{VERLAUF (90 Min):}

\begin{tabular}{|p{1.2cm}|p{5.5cm}|p{5.5cm}|p{1.5cm}|}
\hline
\rowcolor{hellgrau}
\textbf{Zeit} & \textbf{Gruppe A} & \textbf{Gruppe B} & \textbf{L bei} \\
\hline
0--10 & \multicolumn{2}{l|}{Warm-up: Vokabel-Blitzrunde, Fragen klären} & Alle \\
\hline
10--30 & Partnerdialoge üben (wechselnde Partner) & Vorträge üben (2 Durchgänge) & A $\rightarrow$ B \\
\hline
30--50 & Probedurchlauf: 2 Paare präsentieren & Probevortrag vor Gruppe B & B $\rightarrow$ A \\
\hline
50--55 & \multicolumn{2}{l|}{\textit{Pause}} & -- \\
\hline
55--70 & Peer-Feedback mit Kriterienbogen & Peer-Feedback mit Kriterienbogen & A $\rightarrow$ B \\
\hline
70--85 & L-Feedback individuell (kurz!) & L-Feedback individuell (kurz!) & Alle \\
\hline
85--90 & \multicolumn{2}{l|}{Selbsteinschätzung ausfüllen, Abschluss, Motivation} & Alle \\
\hline
\end{tabular}

\vspace{0.2cm}

\textbf{AKTIVITÄTEN:}
\begin{itemize}[leftmargin=1.5em, nosep]
    \item Wechselnde Partnerkonstellationen (jeder übt mit mind. 2 verschiedenen)
    \item Peer-Feedback anhand der Assessment-Kriterien
    \item Selbsteinschätzungsbogen ausfüllen
\end{itemize}

\vspace{0.2cm}

\textbf{INTENTIONEN:}
\begin{itemize}[leftmargin=1.5em, nosep]
    \item Übung unter realistischen Bedingungen (Zeit, Format)
    \item Peer-Feedback fördert Reflexionsfähigkeit
    \item Individuelles L-Feedback gibt letzte Hinweise
    \item \textbf{LETZTE ÜBUNG!} SuS sollen mit Selbstvertrauen ins Assessment gehen
\end{itemize}

\end{dsbox}

\newpage

% ---------- DS 5: ASSESSMENT ----------
\begin{assessmentbox}[title={DS 5 -- ASSESSMENT (Mo, 19.01.2026)}]

\textbf{Thema:} Mündliche Leistungsüberprüfung (Sonstige Note)

\vspace{0.3cm}

\fcolorbox{assessment}{yellow!10}{\parbox{0.95\textwidth}{
\centering
\textbf{\large 45 Minuten | 25 Bewertungseinheiten (BE) | 15-Notenpunkte-Skala (Abitur)}
}}

\vspace{0.3cm}

\textbf{ASSESSMENT-FORMATE:}

\begin{center}
\begin{tabular}{|p{6cm}|p{6cm}|}
\hline
\rowcolor{gruppeA!15}
\textbf{\textcolor{gruppeA}{Gruppe A: Partnerprüfung}} & \textbf{\textcolor{gruppeB}{Gruppe B: Kurzvortrag}} \\
\hline
\textbf{Thema:} \es{Mi barrio} & \textbf{Thema:} \es{Mi profesión ideal} \\
\hline
\textbf{Format:} Dialog (2 SuS, je 2,5 Min) & \textbf{Format:} Einzelvortrag (3 Min + 2 Min Fragen) \\
\hline
\textbf{Dauer:} 5 Min pro Paar & \textbf{Dauer:} 5 Min pro Person \\
\hline
\textbf{Inhalt:} Stadtviertel beschreiben, Fragen stellen/beantworten & \textbf{Inhalt:} Berufswunsch, Begründung, Weg \\
\hline
\end{tabular}
\end{center}

\vspace{0.2cm}

\textbf{VERLAUF (90 Min):}

\begin{tabular}{|p{1.2cm}|p{6cm}|p{5cm}|p{1.5cm}|}
\hline
\rowcolor{hellgrau}
\textbf{Zeit} & \textbf{Aktivität} & \textbf{Parallel} & \textbf{L bei} \\
\hline
0--5 & Begrüßung, Ablauf erklären, Plätze & -- & Alle \\
\hline
5--10 & Vorbereitungszeit (Stichpunkte) & -- & -- \\
\hline
10--20 & \textbf{Gruppe B:} SuS 1 + 2 Vorträge & Gruppe A: Letzte Vorbereitung & B \\
\hline
20--25 & \textbf{Gruppe B:} SuS 3 Vortrag & Gruppe A wartet & B \\
\hline
25--32 & \textbf{Gruppe A:} Paar 1 (E.H./A.S.) & -- & A \\
\hline
32--39 & \textbf{Gruppe A:} Paar 2 (L.N./A.A.) & -- & A \\
\hline
39--46 & \textbf{Gruppe A:} Paar 3 (C.P./V.S.) & -- & A \\
\hline
46--53 & \textbf{Gruppe A:} Paar 4 + Reserve & -- & A \\
\hline
53--60 & Puffer / Nachholprüfungen & Selbstreflexion & Flex \\
\hline
60--75 & Vertiefungsübung (alle gemeinsam) & -- & Alle \\
\hline
75--85 & Kurzes Feedback (Stärken, keine Noten!) & -- & Alle \\
\hline
85--90 & Abschluss, Ausblick DS 6 & -- & Alle \\
\hline
\end{tabular}

\vspace{0.2cm}

\textbf{BEWERTUNGSKRITERIEN (25 BE):}

\begin{center}
\small
\begin{tabular}{|l|c|p{5cm}|p{4.5cm}|}
\hline
\rowcolor{hellgrau}
\textbf{Kriterium} & \textbf{BE} & \textbf{Gruppe A} & \textbf{Gruppe B} \\
\hline
1. Inhalt/Aufgabe & 7,5 & ser/estar/hay korrekt & Argumentation, porque/para \\
\hline
2. Wortschatz & 6,25 & Orte, Adjektive & Berufe, Eigenschaften \\
\hline
3. Grammatik & 5 & Strukturen korrekt & Komplexere Strukturen \\
\hline
4. Flüssigkeit & 6,25 & Verständlich, interaktiv & Frei, gute Aussprache \\
\hline
\multicolumn{2}{|r|}{\textbf{Summe:}} & \multicolumn{2}{l|}{\textbf{25 BE}} \\
\hline
\end{tabular}
\end{center}

\vspace{0.2cm}

\textbf{NOTENSCHLÜSSEL (15-NP-Skala, Abitur):}

\begin{center}
\footnotesize
\begin{tabular}{|c|c|c|c|c|c|c|c|c|c|c|c|c|c|c|c|}
\hline
\textbf{NP} & 15 & 14 & 13 & 12 & 11 & 10 & 9 & 8 & 7 & 6 & 5 & 4 & 3 & 2 & 1 \\
\hline
\textbf{\%} & 98,7 & 97,3 & 96 & 90,7 & 85,3 & 80 & 73,3 & 66,7 & 60 & 53,3 & 46,7 & 40 & 33,3 & 26,7 & 20 \\
\hline
\textbf{BE} & 24,7 & 24,3 & 24 & 22,7 & 21,3 & 20 & 18,3 & 16,7 & 15 & 13,3 & 11,7 & 10 & 8,3 & 6,7 & 5 \\
\hline
\end{tabular}
\end{center}

\vspace{0.2cm}

\textbf{INTENTIONEN:}
\begin{itemize}[leftmargin=1.5em, nosep]
    \item Mündliche Kompetenz als Kernkompetenz im Fremdsprachenunterricht prüfen
    \item Differenzierte Formate (Dialog vs. Monolog) werden beiden Niveaus gerecht
    \item 15-NP-Skala (Abitur) bereitet auf Oberstufe vor
    \item Transparente Kriterien ermöglichen faire Bewertung
\end{itemize}

\end{assessmentbox}

\vspace{0.5cm}

% ═══════════════════════════════════════════════════════════════════════════
% PHASE 2: KONVERGENZ & PUFFER
% ═══════════════════════════════════════════════════════════════════════════

\textbf{\Large\color{spanisch} PHASE 2: Konvergenz \& Puffer (DS 6--8)}

\vspace{0.3cm}

% ---------- DS 6 ----------
\begin{dsbox}[title={DS 6 -- Brückenthema: Personas y características (Di, 20.01.2026)}]

\textbf{Thema:} Erstes vollständig gemeinsames Thema -- Personen beschreiben

\vspace{0.3cm}

\textbf{LERNZIELE:}
\begin{itemize}[leftmargin=1.5em, nosep]
    \item[Alle:] SuS können \es{ser} vs. \es{estar} + Adjektiv kontextabhängig unterscheiden
    \item[Alle:] SuS können Personen mit Charaktereigenschaften beschreiben
    \item[\textcolor{gruppeB}{B:}] SuS können als Tutoren A-SuS bei Schwierigkeiten helfen
\end{itemize}

\vspace{0.2cm}

\textbf{VERLAUF (90 Min):}

\begin{tabular}{|p{1.2cm}|p{9cm}|p{3.5cm}|}
\hline
\rowcolor{hellgrau}
\textbf{Zeit} & \textbf{Aktivität} & \textbf{SF} \\
\hline
0--10 & Begrüßung, Warm-up: Adjektiv-Kette (\es{Soy simpático/a y...}) & Plenum \\
\hline
10--25 & \textbf{Input:} \es{ser} vs. \es{estar} + Adjektiv (Tafel, Beispiele) & Plenum \\
\hline
25--40 & Adjektiv-Sammlung: Brainstorming, Kategorisierung (Charakter/Zustand) & Plenum $\rightarrow$ GA \\
\hline
40--45 & \textit{Pause} & -- \\
\hline
45--65 & \textbf{Gemischte Tandems:} \es{¿Cómo eres?} (B interviewt A) & PA (A+B) \\
\hline
65--80 & Spiel: \es{¿Quién es?} (Personen erraten anhand Beschreibung) & Plenum \\
\hline
80--90 & Abschluss, Reflexion, HA: Lieblingsperson beschreiben (5 Sätze) & Plenum \\
\hline
\end{tabular}

\vspace{0.2cm}

\textbf{AKTIVITÄTEN:}
\begin{itemize}[leftmargin=1.5em, nosep]
    \item Grammatik-Input mit kontrastiven Beispielen (ser = permanent, estar = temporär)
    \item Gemischte Tandems: B-SuS als Tutoren, A-SuS profitieren von Sprachmodellen
    \item Spielerische Anwendung: \es{¿Quién es?} fördert Hörverstehen und Sprechfreude
\end{itemize}

\vspace{0.2cm}

\textbf{INTENTIONEN:}
\begin{itemize}[leftmargin=1.5em, nosep]
    \item Erstes \textbf{gemeinsames} Thema nach paralleler Bucharbeit stärkt Gemeinschaft
    \item \textcolor{gruppeB}{\textbf{B-SuS als Tutoren}} entlastet Lehrkraft und fördert B-Selbstbewusstsein
    \item Vorbereitung auf Beobachtungsstunde (DS 10): Tandem-Format wird eingeübt
    \item \es{ser/estar} als gemeinsame grammatische Grundlage für alle weiteren Stunden
\end{itemize}

\end{dsbox}

\vspace{0.5cm}

% ---------- DS 7: PUFFER ----------
\begin{pufferbox}[title={DS 7 -- PUFFER: Cultura y Diversión (Mo, 26.01.2026)}]

\textbf{Thema:} Kulturelle Stunde ohne neuen prüfungsrelevanten Stoff

\vspace{0.3cm}

\fcolorbox{puffer}{purple!5}{\parbox{0.95\textwidth}{
\centering
\textbf{PUFFERWOCHE!} Kein neuer Stoff. Entbehrlich für Beobachtung.
}}

\vspace{0.3cm}

\textbf{LERNZIELE:}
\begin{itemize}[leftmargin=1.5em, nosep]
    \item[Alle:] SuS erweitern ihr kulturelles Wissen über spanischsprachige Länder
    \item[Alle:] SuS wenden bekannten Wortschatz spielerisch an
    \item[Alle:] SuS erleben Spanisch als lebendige Sprache (Motivation)
\end{itemize}

\vspace{0.2cm}

\textbf{VERLAUF (90 Min):}

\begin{tabular}{|p{1.2cm}|p{9cm}|p{3.5cm}|}
\hline
\rowcolor{hellgrau}
\textbf{Zeit} & \textbf{Aktivität} & \textbf{SF} \\
\hline
0--10 & Begrüßung, Warm-up: \es{¿Qué sabes de España/Latinoamérica?} & Plenum \\
\hline
10--30 & Landeskundequiz (Kahoot oder analog) -- Länder, Hauptstädte, Fakten & Teams \\
\hline
30--45 & Spanische Musik hören, Lückentext ausfüllen (einfaches Lied) & EA $\rightarrow$ PA \\
\hline
45--50 & \textit{Pause} & -- \\
\hline
50--70 & Vokabel-Bingo (Stadtviertel + Eigenschaften + Berufe) & Plenum \\
\hline
70--85 & Kurzfilm (5 Min) mit einfachen Verständnisfragen & Plenum \\
\hline
85--90 & Abschluss, Ausblick auf DS 8 & Plenum \\
\hline
\end{tabular}

\vspace{0.2cm}

\textbf{AKTIVITÄTEN:}
\begin{itemize}[leftmargin=1.5em, nosep]
    \item Quiz: Wettbewerb in Teams (A+B gemischt)
    \item Musik: Eingängiges spanisches Lied (z.B. \es{Bailando}, \es{Despacito} -- didaktisierte Version)
    \item Bingo: Wiederholung aller wichtigen Vokabeln
    \item Kurzfilm: Authentisches Material, einfache Aufgaben
\end{itemize}

\vspace{0.2cm}

\textbf{INTENTIONEN:}
\begin{itemize}[leftmargin=1.5em, nosep]
    \item \textbf{Puffer} falls Assessment länger dauert oder Nachholtermine nötig sind
    \item Motivationsstunde nach Assessment -- Belohnung für gute Vorbereitung
    \item Kulturelle Kompetenz als oft vernachlässigter Kompetenzbereich
    \item Online-Variante verfügbar bei Lehrkraft-Abwesenheit
\end{itemize}

\end{pufferbox}

\vspace{0.5cm}

% ---------- DS 8: PUFFER ----------
\begin{pufferbox}[title={DS 8 -- PUFFER: Creatividad y Juego (Di, 27.01.2026)}]

\textbf{Thema:} Kreative Anwendung -- \es{Mi persona favorita}

\vspace{0.3cm}

\fcolorbox{puffer}{purple!5}{\parbox{0.95\textwidth}{
\centering
\textbf{PUFFERWOCHE!} Kein neuer Stoff. Entbehrlich für Beobachtung.
}}

\vspace{0.3cm}

\textbf{LERNZIELE:}
\begin{itemize}[leftmargin=1.5em, nosep]
    \item[Alle:] SuS können eine Lieblingsperson schriftlich beschreiben
    \item[\textcolor{gruppeA}{A:}] SuS schreiben mind. 60--80 Wörter mit bekannten Strukturen
    \item[\textcolor{gruppeB}{B:}] SuS schreiben 100--120 Wörter mit komplexeren Strukturen
\end{itemize}

\vspace{0.2cm}

\textbf{VERLAUF (90 Min):}

\begin{tabular}{|p{1.2cm}|p{9cm}|p{3.5cm}|}
\hline
\rowcolor{hellgrau}
\textbf{Zeit} & \textbf{Aktivität} & \textbf{SF} \\
\hline
0--10 & Begrüßung, Brainstorming: \es{¿Quién es tu persona favorita?} & Plenum \\
\hline
10--20 & Wortschatz-Wiederholung: Eigenschaften, Aussehen, Aktivitäten & Plenum \\
\hline
20--45 & \textbf{Schreibphase:} \es{Mi persona favorita} (differenziert) & EA \\
\hline
45--50 & \textit{Pause} & -- \\
\hline
50--65 & Freiwillige Präsentation (2--3 Texte vorlesen) & Plenum \\
\hline
65--80 & Padlet/Poster: Texte digital oder analog teilen & GA \\
\hline
80--90 & Feedback-Runde, Abschluss & Plenum \\
\hline
\end{tabular}

\vspace{0.2cm}

\textbf{AKTIVITÄTEN:}
\begin{itemize}[leftmargin=1.5em, nosep]
    \item Schreibaufgabe mit Scaffolding (Satzanfänge, Wortschatzlisten)
    \item Freiwillige Präsentation (kein Zwang!)
    \item Digitale oder analoge Sammlung der Texte
\end{itemize}

\vspace{0.2cm}

\textbf{INTENTIONEN:}
\begin{itemize}[leftmargin=1.5em, nosep]
    \item Schreibkompetenz als Vorbereitung auf schriftliche Prüfungen
    \item Persönlicher Bezug (\es{persona favorita}) erhöht Motivation
    \item Keine neue Grammatik -- nur Anwendung bekannter Strukturen
    \item Letzter Puffer vor der Beobachtungswoche
\end{itemize}

\end{pufferbox}

\newpage

% ═══════════════════════════════════════════════════════════════════════════
% PHASE 3: EINSTIMMUNG & BEOBACHTUNG
% ═══════════════════════════════════════════════════════════════════════════

\textbf{\Large\color{spanisch} PHASE 3: Einstimmung \& Beobachtung (DS 9--10)}

\vspace{0.3cm}

% ---------- DS 9 ----------
\begin{dsbox}[title={DS 9 -- Repaso: Vocabulario y Gramática (Mo, 02.02.2026)}]

\textbf{Thema:} Spielerische Wiederholung -- Vokabeln und Grammatik festigen

\vspace{0.3cm}

\fcolorbox{spanisch}{hellgrau}{\parbox{0.95\textwidth}{
\centering
\textbf{WICHTIG:} Komplett anderer Ansatz als DS 10! \\
Keine Tandem-Interviews, keine \glqq{}Generalprobe\grqq{} -- SuS sollen nichts doppelt machen.
}}

\vspace{0.3cm}

\textbf{LERNZIELE:}
\begin{itemize}[leftmargin=1.5em, nosep]
    \item[Alle:] SuS können \es{ser/estar/hay} sicher anwenden
    \item[Alle:] SuS beherrschen Vokabular (Eigenschaften, Berufe, Hobbys)
    \item[Alle:] SuS fühlen sich sicher und selbstbewusst
\end{itemize}

\vspace{0.2cm}

\textbf{VERLAUF (90 Min):}

\begin{tabular}{|p{1.2cm}|p{9cm}|p{3.5cm}|}
\hline
\rowcolor{hellgrau}
\textbf{Zeit} & \textbf{Aktivität} & \textbf{SF} \\
\hline
0--10 & Begrüßung, Warm-up: Adjektiv-Kette (Ich bin... und...) & Plenum \\
\hline
10--30 & \textbf{Vokabel-Bingo:} Eigenschaften + Berufe (spielerisch!) & GA \\
\hline
30--45 & \textbf{Grammatik-Stationen:} ser vs. estar Übungen (3 Stationen) & Stationen \\
\hline
45--50 & \textit{Pause} & -- \\
\hline
50--70 & \textbf{Schreibübung:} Kurzer Text über Lieblingsstar (3. Person) & EA \\
\hline
70--80 & Texte vorlesen lassen, Peer-Feedback & Plenum \\
\hline
80--85 & \textbf{Kurze Ankündigung:} Format für morgen verbal erklären (5 Min) & Plenum \\
\hline
85--90 & Motivierende Worte, offene Fragen, Abschluss & Plenum \\
\hline
\end{tabular}

\vspace{0.2cm}

\textbf{AKTIVITÄTEN:}
\begin{itemize}[leftmargin=1.5em, nosep]
    \item Vokabel-Bingo: Spielerisch, ohne Druck
    \item Grammatik-Stationen: Bewegung im Raum, abwechslungsreich
    \item Schreibübung: Festigt 3. Person (Vorbereitung ohne Probe)
    \item \textbf{Nur 5 Min verbal:} Format morgen kurz erklären (keine Demonstration!)
\end{itemize}

\vspace{0.2cm}

\textbf{INTENTIONEN:}
\begin{itemize}[leftmargin=1.5em, nosep]
    \item \textbf{Kein Probedurchlauf!} DS 10 soll frisch wirken, nicht künstlich
    \item Positive Atmosphäre: Spiele reduzieren Stress
    \item Schreibübung übt 3. Person ohne Interview-Situation
    \item Morgen \glqq{}Format\grqq{} nur verbal ankündigen, nicht demonstrieren
\end{itemize}

\end{dsbox}

\vspace{0.5cm}

% ---------- DS 10: BEOBACHTUNG ----------
\begin{beobachtungsbox}[title={DS 10 -- BEOBACHTUNGSSTUNDE (Di, 03.02.2026)}]

\textbf{Thema:} Tandem-Interview \es{¿Cómo eres?} -- Externe Unterrichtsbeobachtung

\vspace{0.3cm}

\fcolorbox{beobachtung}{red!5}{\parbox{0.95\textwidth}{
\centering
\textbf{\large EXTERNE BEOBACHTUNG!}\\[0.2cm]
Beobachtet wird die \textbf{2. Hälfte} (45 Min)\\[0.1cm]
Die 1. Hälfte dient der Vorbereitung
}}

\vspace{0.3cm}

\textbf{LERNZIELE:}
\begin{itemize}[leftmargin=1.5em, nosep]
    \item[Alle:] SuS führen ein strukturiertes Interview auf Spanisch
    \item[Alle:] SuS verwenden \es{ser/estar/gustar} korrekt im Kontext
    \item[\textcolor{gruppeB}{B:}] SuS fassen Informationen in der 3. Person Singular zusammen
    \item[\textcolor{gruppeB}{B:}] SuS unterstützen A-SuS als Peer-Tutoren
\end{itemize}

\vspace{0.2cm}

\textbf{VERLAUF 1. HÄLFTE (VOR Beobachtung -- 45 Min):}

\begin{tabular}{|p{1.2cm}|p{9cm}|p{3.5cm}|}
\hline
\rowcolor{hellgrau}
\textbf{Zeit} & \textbf{Aktivität} & \textbf{SF} \\
\hline
0--10 & Begrüßung, Warm-up: Adjektiv-Kette & Plenum \\
\hline
10--25 & Wiederholung in homogenen Gruppen (A mit A, B mit B) & GA \\
\hline
25--40 & Interview-Fragen durchgehen, Regeln erklären, Beispiel & Plenum \\
\hline
40--45 & Pause, Beobachterin kommt an & -- \\
\hline
\end{tabular}

\vspace{0.2cm}

\textbf{VERLAUF 2. HÄLFTE (BEOBACHTUNG -- 45 Min):}

\begin{tabular}{|p{1.2cm}|p{9cm}|p{3.5cm}|}
\hline
\rowcolor{red!10}
\textbf{Zeit} & \textbf{Aktivität} & \textbf{SF} \\
\hline
0--3 & Orientierung: Ablauf, Timer, Regeln (\es{Solo español!}) & Plenum \\
\hline
3--15 & \textbf{Phase 1:} Homogene Interviews (A$\leftrightarrow$A, B$\leftrightarrow$B) -- 12 Min & PA \\
\hline
15--33 & \textbf{Phase 2:} Gemischte Tandems -- B interviewt A, \textbf{Ankreuz-Vorlage} -- 18 Min & PA (A+B) \\
\hline
33--42 & \textbf{Phase 3:} Präsentation -- \textbf{nur 3 Tandems}, B stellt A vor -- 9 Min & Plenum \\
\hline
42--45 & Cierre: Zusammenfassung, Verabschiedung auf Spanisch & Plenum \\
\hline
\end{tabular}

\vspace{0.2cm}

\fcolorbox{success}{success!10}{\parbox{0.95\textwidth}{
\textbf{Timing-Optimierung:} Ankreuz-Vorlage statt freiem Notieren = schneller. Nur 3 vorher festgelegte Tandems präsentieren = keine Zeitüberschreitung.
}}

\vspace{0.2cm}

\textbf{DIFFERENZIERTE FRAGEKARTEN:}

\begin{center}
\begin{tabular}{|p{6cm}|p{6cm}|}
\hline
\rowcolor{gruppeA!15}
\textbf{\textcolor{gruppeA}{Fragekarten A (3 Fragen)}} & \textbf{\textcolor{gruppeB}{Fragekarten B (4 Fragen)}} \\
\hline
\es{¿Cómo te llamas?} & \es{¿Qué profesión quieres tener?} \\
\es{¿Cómo eres?} (+ Beispiele) & \es{¿Qué cualidades son importantes?} \\
\es{¿Qué te gusta?} (+ Beispiele) & \es{¿Qué actividades haces?} \\
 & \es{¿Cómo te describirías?} \\
\hline
\end{tabular}
\end{center}

\vspace{0.2cm}

\textbf{TANDEM-ZUORDNUNG (B interviewt A):}

\begin{center}
\begin{tabular}{|c|l|l|l|}
\hline
\textbf{Nr.} & \textbf{Gruppe B} & \textbf{Gruppe A} & \textbf{Begründung} \\
\hline
1 & E.H. (TOP) & H.P. (Scaffolding) & Stärkste/r unterstützt Schwächste/n \\
\hline
2 & L.N. (proaktiv) & A.A. (unsicher) & Ermutigung durch aktive/n Partner/in \\
\hline
3 & V.S. (konstant) & A.S. (einziger Junge) & Ruhige, unterstützende Atmosphäre \\
\hline
4 & C.P. (ruhig) & L.K. (solide) & Ähnliches Tempo \\
\hline
5 & [S9] & [S10] & -- \\
\hline
\end{tabular}
\end{center}

\vspace{0.2cm}

\textbf{INTENTIONEN (Was die Beobachterin sehen soll):}
\begin{itemize}[leftmargin=1.5em, nosep]
    \item[$\checkmark$] \textbf{Differenzierung:} Zwei Komplexitätsniveaus, Peer-Scaffolding
    \item[$\checkmark$] \textbf{Aktivierung:} Alle SuS sprechen mehrfach, kein passives Zuhören
    \item[$\checkmark$] \textbf{Kommunikation:} Echte Interaktion, nicht nur Ablesen
    \item[$\checkmark$] \textbf{Scaffolding:} Fragekarten, Strukturhilfen, B als Tutoren
    \item[$\checkmark$] \textbf{Zielsprachenverwendung:} Maximaler Spanisch-Anteil
\end{itemize}

\vspace{0.3cm}

\textbf{ALTERNATIVE ENTWÜRFE (als Backup):}

\begin{center}
\begin{tabular}{|l|p{4cm}|c|c|}
\hline
\rowcolor{hellgrau}
\textbf{Variante} & \textbf{Ansatz} & \textbf{Timing-Risiko} & \textbf{Score} \\
\hline
\textbf{A (Standard)} & Tandem-Interview + Ankreuz-Vorlage & Mittel & 9.4/10 \\
\hline
\textbf{B} & Speed-Dating (6×3 Min, kein Notieren) & \textcolor{success}{Niedrig} & 8.7/10 \\
\hline
\textbf{C} & Galerierundgang (Poster + Gallery Walk) & \textcolor{success}{Niedrig} & 9.0/10 \\
\hline
\end{tabular}
\end{center}

{\small Empfehlung: Variante A zeigt beste Differenzierung. B/C als Backup bei Unsicherheit.}

\end{beobachtungsbox}

\newpage

% ═══════════════════════════════════════════════════════════════════════════
% GRAMMATISCHE PROGRESSION
% ═══════════════════════════════════════════════════════════════════════════

{\large\bfseries\color{spanisch} Grammatische Progression}

\vspace{0.3cm}

\begin{center}
\begin{tabular}{|p{3cm}|p{5.5cm}|p{5.5cm}|}
\hline
\rowcolor{hellgrau}
\textbf{Struktur} & \textbf{\textcolor{gruppeA}{Gruppe A (A2)}} & \textbf{\textcolor{gruppeB}{Gruppe B (B1)}} \\
\hline
\es{ser} + Adjektiv & DS 2: Einführung (einfache Eigenschaften) & DS 1: Wiederholung (bekannt) \\
\hline
\es{estar} + Adjektiv & DS 6: Einführung (Zustände) & DS 1: Wiederholung (bekannt) \\
\hline
\es{hay} + Ortsangabe & DS 1: Einführung (\es{Mi barrio}) & -- (bereits bekannt) \\
\hline
\es{gustar} & DS 2: Wiederholung (Vorlieben) & DS 1: Anwendung (Hobbys) \\
\hline
\es{querer ser} & -- & DS 1--5: Fokus (Berufswunsch) \\
\hline
3. Person Singular & DS 6: Einführung (\es{él/ella es...}) & DS 6+10: Anwendung (Präsentation) \\
\hline
\es{llamarse} & DS 1: bekannt & DS 1: bekannt \\
\hline
Kausale Konj. & -- & DS 1--5: \es{porque, ya que, para} \\
\hline
\end{tabular}
\end{center}

\vspace{0.5cm}

% ═══════════════════════════════════════════════════════════════════════════
% MATERIALIEN
% ═══════════════════════════════════════════════════════════════════════════

{\large\bfseries\color{spanisch} Verfügbare Materialien}

\vspace{0.3cm}

\begin{tabular}{@{}p{4cm}p{5cm}p{6cm}@{}}
\textbf{Typ} & \textbf{Dateien} & \textbf{Beschreibung} \\
\hline
Langentwürfe & LE-01 bis LE-10 (.tex/.pdf) & Je ca. 8--12 Seiten, Tier 1 (vollständig) \\
Fragekarten & MATERIAL-10-fragekarten.pdf & 6x A (3 Fragen), 6x B (4 Fragen), Poster \\
Bewertungsbögen & MATERIAL-05-bewertungsboegen.pdf & Assessment: 25 BE, 15 NP \\
Zusammenfassung & Resumen-Vorlage (Gruppe B) & Für Präsentation DS 10 \\
Materialübersicht & UEBERSICHT-materialien.pdf & Komplette Materialliste mit Checkliste \\
Diese Übersicht & SEQUENZ-uebersicht.pdf & Kompakte Stundenübersicht \\
\end{tabular}

\vspace{0.5cm}

% ═══════════════════════════════════════════════════════════════════════════
% CHECKLISTE
% ═══════════════════════════════════════════════════════════════════════════

{\large\bfseries\color{spanisch} Checkliste vor der Beobachtung}

\vspace{0.3cm}

\begin{tabular}{|p{0.5cm}|p{14cm}|}
\hline
$\Box$ & Alle 10 Langentwürfe ausgedruckt und gelesen \\
\hline
$\Box$ & Fragekarten A und B laminiert (je 6 Stück) \\
\hline
$\Box$ & Poster \es{ser vs. estar} vorbereitet \\
\hline
$\Box$ & Tandem-Übersicht für Lehrkraft bereit (5 Paare) \\
\hline
$\Box$ & Timer/Uhr sichtbar eingerichtet \\
\hline
$\Box$ & Bewertungsbögen für Assessment (DS 5) gedruckt -- \textbf{25 BE, 15 NP} \\
\hline
$\Box$ & Sitzordnung für gemischte Tandems geplant \\
\hline
$\Box$ & DS 9 durchgeführt (Vokabel-/Grammatikspiele, KEIN Probedurchlauf) \\
\hline
$\Box$ & Zusammenfassungs-Vorlage für Gruppe B kopiert \\
\hline
$\Box$ & Strukturhilfe (1. $\rightarrow$ 3. Person) an Tafel vorbereitet \\
\hline
$\Box$ & Beobachterin über Ablauf informiert (2. Hälfte!) \\
\hline
\end{tabular}

\vspace{0.5cm}

% ═══════════════════════════════════════════════════════════════════════════
% NOTIZEN
% ═══════════════════════════════════════════════════════════════════════════

{\large\bfseries\color{spanisch} Notizen}

\vspace{0.3cm}

\begin{tabular}{|p{15.5cm}|}
\hline
\\[5cm]
\hline
\end{tabular}

\end{document}
